\subsection{模的n次外幂与交错多重线性映射}

给定交换环 A 上的两个模 M、N，从 \(M^{n}\) 到 N 的交错 n-线性映射是指这样的 n-线性映射 \(f: M^{n} \to N\) ：对所有 \(p \leqslant n - 2\) ，

\[
f (u _ {1}, \dots , u _ {p}, x, x, v _ {1}, \dots , v _ {n - p - 2}) = 0\tag{8}
\]

对所有 \(x\) ，所有 \(u_{i}\) ( \(1 \leqslant i \leqslant p\) ) 以及所有 \(v_{j}\) ( \(1 \leqslant j \leqslant n - p - 2\) ) 在 \(\mathbf{M}\) 中均成立。

命题7. 设 A 为交换环，M 和 N 为两个 A-模。若对每个 A-线性映射 \(g: \bigwedge^n (\mathbf{M}) \to \mathbf{N} (n \geqslant 2)\) 使之对应于 n-线性映射

\[
\left(x _ {1}, x _ {2}, \dots , x _ {n}\right) \mapsto g \left(x _ {1} \wedge x _ {2} \wedge \dots \wedge x _ {n}\right)\tag{9}
\]

便得到从 A-模 \(\operatorname{Hom}_{\mathbf{A}}(\bigwedge^{n}(\mathbf{M}), \mathbf{N})\) 到由交错 \(n\) -线性映射（从 \(\mathbf{M}^n\) 到 \(\mathbf{N}\) 的）所组成的 A-模上的一个双射 A-线性映射。

我们考虑 A-模 \(\operatorname{Hom}_{\mathbf{A}}(\mathsf{T}^{n}(\mathbf{M}),\mathbf{N})\) 到 A-模 \(\mathcal{L}_{n}(\mathbf{M},\ldots,\mathbf{M};\mathbf{N})\) （即所有从 \(M^{n}\) 到 N 的 n-线性映射组成的 A-模）的典范双射，它通过将每个 A-线性映射 \(f: T^{n}(\mathbf{M}) \to \mathbf{N}\) 对应于 n-线性映射而得到

\[
\bar {f} \colon (x _ {1}, \dots , x _ {n}) \mapsto f (x _ {1} \otimes x _ {2} \otimes \dots \otimes x _ {n})
\]

(II, § 3, no. 9)。另一方面，A-线性映射 \(g \colon \bigwedge^{n}(\mathbf{M}) \to \mathbf{N}\) 与 A-线性映射 \(f \colon \mathsf{T}^n (\mathbf{M}) \to \mathbf{N}\) （满足 \(f\) 在 \(\mathfrak{J}_n''\) 上为零者）一一对应，其对应方式是把 \(g\) 映为 \(f = g \circ p_n\) ，其中

\[
p _ {n} \colon \mathsf {T} ^ {n} (\mathbf {M}) \rightarrow \bigwedge^ {n} (\mathbf {M}) = \mathsf {T} ^ {n} (\mathbf {M}) / \mathfrak {J} _ {n} ^ {\prime \prime}
\]

是典范同态（II，§2，第 1 号，定理 1）。但鉴于 \(\mathfrak{J}_n''\) 是形如

\[
\left(u _ {1} \otimes u _ {2} \otimes \dots \otimes u _ {p}\right) \otimes (x \otimes x) \otimes \left(v _ {1} \otimes \dots \otimes v _ {n - p - 2}\right)
\]

\((x, u_{i}, v_{j} \text{ 属于 M}), \text{ 说 } f \text{ 具有形式 } g \circ p_{n} \text{ 意味着相应的 } n\text{-线性函数 } \bar{f} \text{ 满足 (8)，换言之它是交错的。}\)

A-模 \(\bigwedge^{n}(\mathbf{M})\) 称为 \(n\) -次外幂（即 \(\mathbf{M}\) 的外幂）。对于每个 A-模同态 \(u: \mathbf{M} \to \mathbf{N}\) ，映射

\[
\bigwedge^ {n} (u): \bigwedge^ {n} (\mathbf {M}) \rightarrow \bigwedge^ {n} (\mathbf {N})
\]

它与 \(\bigwedge(u)\) 在 \(\bigwedge^{n}(\mathbf{M})\) 上的限制相一致，称为 \(n\) -次外幂（即 \(u\) 的外幂）.

推论 1. 对于每个交错 \(n\) -线性映射 \(g: \mathbf{M}^n \to \mathbf{N}\) ，对于每个置换 \(\sigma \in \mathfrak{S}_n\) ,

\[
g \left(x _ {\sigma (1)}, x _ {\sigma (2)}, \dots , x _ {\sigma (n)}\right) = \varepsilon_ {\sigma}. g \left(x _ {1}, x _ {2}, \dots , x _ {n}\right)\tag{10}
\]

对所有 \(x_{i} \in \mathbf{M}\) 成立；此外，若 \(x_{i} = x_{j}\) 对于两个不同的指标 \(i, j\) 成立，则

\[
g (x _ {1}, x _ {2}, \dots , x _ {n}) = 0.
\]

这是命题 7 以及第 3 号命题 5 的明显推论。

推论2. 设 \((x_{i})_{1\leqslant i\leqslant n}\) 是由 \(n\) 个元素组成的一个序列（元素属于 \(\mathbf{M}\) ），使得

\[
x _ {1} \wedge x _ {2} \wedge \dots \wedge x _ {n} = 0;
\]

那么，对于每一个交错的 \(n\) -线性映射 \(g: \mathbf{M}^n \to \mathbf{N}\) , \(g(x_1, \ldots, x_n) = 0\) .

推论3. 设 \(f: \mathbf{M}^{n-1} \to \mathbf{A}\) 是一个交错 \((n-1)\) -线性形式。如果 \((x_i)_{1 \leq i \leq n}\) 是由 \(n\) 个元素组成的一个序列（元素属于 \(\mathbf{M}\) ），使得 \(x_1 \wedge x_2 \wedge \cdots \wedge x_n = 0\) ，则

\[
\sum_ {i = 1} ^ {n} (- 1) ^ {i} f (x _ {1}, \dots , \hat {x _ {i}}, \dots , x _ {n}). x _ {i} = 0\tag{11}
\]

（其中我们记 \(f(x_{1},\ldots ,\hat{x}_{i},\ldots ,x_{n}) = f(x_{1},\ldots ,x_{i - 1},x_{i + 1},\ldots ,x_{n})\) 对于 \(1\leqslant i\leqslant n\) .

只需证明这个 n-线性映射

\[
(x _ {1}, \dots , x _ {n}) \mapsto \sum_ {i = 1} ^ {n} (- 1) ^ {i} f (x _ {1}, \dots , \hat {x} _ {i}, \dots , x _ {n}). x _ {i}
\]

从 \(M^{n}\) 到 M 是交错的。现在，若 \(x_{i} = x_{i+1}\) ，则由于 f 是交错的，右端求和中除 \(x_{i}\) 与 \(x_{i+1}\) 外的所有项系数均为零；另一方面， \(x_{i}\) 的系数是

\[
(- 1) ^ {i} f (x _ {1}, \dots , x _ {i - 1}, x _ {i + 1}, x _ {i + 2}, \dots , x _ {n})
\]

而 \(x_{i+1}\) 的系数是 \((-1)^{i+1}f(x_1, \ldots, x_i, x_{i+2}, \ldots, x_n)\) ，由假设二者互为相反数。

注记。一个元素 \(z\) 在 \(\mathsf{T}^n (\mathbf{M})\) 中称为 \(n\) 阶（逆变）反对称张量，若 \(\sigma .z = \varepsilon_{\sigma}z\) （对每个置换 \(\sigma \in \mathfrak{S}_n\) ；参见 §6，第 3 号，注记）；这些元素构成一个子 A-模 \(\mathsf{A}_n^\prime (\mathbf{M})\) ，它包含于 \(\mathsf{T}^n (\mathbf{M})\) 。对于所有 \(z\in \mathsf{T}^n (\mathbf{M})\) ，我们记 \(\pmb {a}.z = \sum_{\sigma \in \mathfrak{S}_n}\varepsilon_\sigma (\sigma .z)\) ，并称 \(\pmb {a}.z\) 为 \(z\) 的反对称化；由于 \(\varepsilon_{\sigma \tau} = \varepsilon_{\sigma}\varepsilon_{\tau}\) ，立即可见 \(\pmb {a}.z\) 是反对称张量，因而 \(z\mapsto \pmb {a}.z\) 是 \(\mathsf{T}^n (\mathbf{M})\) 的一个自同态，其像 \(\mathsf{A}_n''(\mathbf{M})\) 包含于 \(\mathsf{A}_n^\prime (\mathbf{M})\) ；一般地 \(\mathsf{A}_n''(\mathbf{M})\neq \mathsf{A}_n^\prime (\mathbf{M})\) （习题 8）。若 \(z\) 是反对称张量，则 \(\pmb {a}.z = n!z\) ；因此，当 \(n!\) 在 A 中可逆时，自同态 \(z\mapsto (n!)^{-1}\pmb {a}.z\) 是 \(\mathsf{T}^n (\mathbf{M})\) 的一个投影算子，其像为

\[
\mathbf {A} _ {n} ^ {\prime} (\mathbf {M}) = \mathbf {A} _ {n} ^ {\prime \prime} (\mathbf {M}).
\]

此外，这个投影算子的核正是 \(\mathfrak{J}_n''\) ；因为显然只需限于 \(n \geqslant 2\) 的情形，于是 2（它是 \(n!\) 的因子）在 A 中可逆，且 \(x \otimes x = 2^{-1}(x \otimes x + x \otimes x)\) ；因此 \(\mathfrak{J}_n''\) 由元素 \(z + \rho . z\) 生成，其中 \(\rho\) 是交换 \([1, n]\) 中两个相邻数的对换；此外，对于两个置换 \(\sigma, \tau\) （在 \(\mathfrak{S}_n\) 中），我们可以写成

\[
z - \varepsilon_ {\sigma \tau} ((\sigma \tau). z) = z - \varepsilon_ {\tau} (\tau . z) + \varepsilon_ {\tau} (\tau . z - \varepsilon_ {\sigma} \sigma . (\tau . z))
\]

由此可得 \(z - \varepsilon_{\sigma}(\sigma . z) \in \mathfrak{J}_{n}^{\prime \prime}\) 对所有 \(z \in \mathbf{T}^{n}(\mathbf{M})\) 和 \(\sigma \in \mathfrak{S}_{n}\) 成立。因此（始终假设 \(n!\) 在 A 中可逆），可以看出

\[
z - (n!) ^ {- 1} \boldsymbol {a}. z = \sum_ {\sigma \in \mathfrak {S} _ {n}} (n!) ^ {- 1} (z - \varepsilon_ {\sigma} (\sigma . z)) \in \mathfrak {J} _ {n} ^ {\prime \prime}
\]

对所有 \(z \in \mathbf{T}^n(\mathbf{M})\) 成立，这确立了我们的断言。

当 \(n!\) 在 A 中可逆时，子模 \(\mathsf{A}_n'(M)\) 和 \(\mathfrak{J}_n''\) （均为 \(T^n(M)\) 的子模）因此是互补的，并且限制到 \(\mathsf{A}_n'(M)\) 上的典范同态 \(T^n(M) \to \bigwedge^n(M) = T^n(M)/\mathfrak{J}_n''\) 是一个 A-模同构，这使我们在所讨论的情形下可以把 \(n\) 阶反对称张量与 M 的 \(n\) -次外幂的元素等同起来。这里还需注意，这一等同与乘法不相容：两个反对称张量在 \(T(M)\) 中的乘积一般不再是反对称的。

